\section{The Shaders: Same Game, Different Bytes Per Vendor}
\label{sec:shaders}

The driver crashes compiling a shader. The obvious next question is \emph{which}
shader, and the obvious assumption is that both GPUs receive the same shaders and
only one vendor's compiler mishandles them. The second half of that assumption is
false, and correcting it reframes the problem.

\subsection{Dump both, and compare}

The translation layer can be told to write every shader it produces, in both its
input and output forms, to a directory. Running the game once on each GPU and
comparing the two directories gives the following.

\begin{table}[htbp]
\caption{Shader output compared across the two GPUs, same game, same build}
\label{tab:shaders}
\centering
\footnotesize
\begin{tabularx}{\columnwidth}{@{}L{4.4cm}X@{}}
\toprule
\textbf{Measurement} & \textbf{Value} \\
\midrule
Modules translated, Radeon & 6{,}657 \\
Modules translated, NVIDIA & 4{,}430 \\
Hashes present in both sets & 4{,}430 \\
Same hash, \textbf{different bytes} & \textbf{4{,}047} \\
Uses \verb|SPV_NV_raw_access_chains| & 4{,}039 NVIDIA / \textbf{0} Radeon \\
\quad of which compute shaders & 561 of 886 \\
Uses \verb|SPV_NV_shader_|\linebreak\verb|subgroup_partitioned| & exactly 1
                                                                 compute shader \\
\bottomrule
\end{tabularx}
\end{table}

The interpretation is unambiguous. \textbf{The translation layer emits
vendor-specific code.} Given identical Direct3D shader bytecode as input, it
produces different Vulkan output depending on which driver is present, because it
uses vendor extensions where they are available.

Of the shaders present on both sides, \textbf{91\% differ byte-for-byte}. The
NVIDIA path uses a vendor extension in 4{,}039 modules that the Radeon path uses
in none. So the premise ``the same shaders work on AMD'' was never true. AMD was
never given these shaders.

This matters for a reason beyond diagnosis: it means the AMD success provided
much weaker evidence than it appeared to. It did not show that the game's shaders
are sound; it showed that a \emph{different} set of shaders is sound.

\subsection{Overriding the shaders, and a large partial result}

The translation layer also supports substituting a shader by hash: place a
replacement in a directory and the corresponding module is loaded from there
instead of being generated. Feeding the Radeon's modules to the NVIDIA run is
therefore mechanically possible.

The result was the single largest improvement of the investigation:

\begin{center}
\texttt{9 dialogs, GPU 0\%} $\longrightarrow$
\textbf{\texttt{4 dialogs, GPU 100\%}}
\end{center}

Five of the nine crashing pipelines stopped crashing, and the discrete GPU went
from idle to fully loaded---it was doing real work for the first time. The result
held across \textbf{fourteen consecutive invariant runs}, so it is not noise.

It is also not a fix. Four dialogs remain, and the configuration is not one that
could ship: it requires a shader dump from the other GPU, and it forces the
translation layer to bypass its own pipeline cache.

\subsection{Why this experiment is not clean, stated plainly}

It is tempting to read the partial success as ``the vendor extension is the
bug.'' The experiment does not support that, for two reasons that must be stated:

\begin{enumerate}
\item Substituting shaders \textbf{forces the layer to ignore its cached SPIR-V
      entirely}, changing a second variable at the same time as the first.
\item The Radeon modules were compiled for a different device's capability set.
      Feeding them to a different vendor's compiler is not a controlled
      substitution; it is a different program.
\end{enumerate}

\subsection{An unresolved contradiction, published rather than smoothed over}

If the vendor extension were the cause, then disabling that extension outright
should reproduce the improvement. It does not.

\begin{quote}
Disabling \verb|VK_NV_raw_access_chains| through the layer's own extension
filter---so that the same 4{,}039 modules are generated without it---does
\textbf{not} reduce the dialog count. The override, which changes both the
extension usage and the cache path, does.
\end{quote}

We were unable to resolve this contradiction, and we report it unresolved. The
most likely explanations are that the improvement came from the cache bypass
rather than the extension, or that the extension filter does not affect
already-cached modules. Neither was confirmed. It is recorded here because a
reader reproducing this work will hit the same contradiction, and because the
correct fix, found later, makes the question moot in a way that is itself
informative: the layer that emitted the shaders was the wrong layer to be
adjusting.

\subsection{Everything tried at the shader level}

\begin{table}[htbp]
\caption{Shader- and translation-layer configuration attempts, all negative}
\label{tab:shaderelim}
\centering
\footnotesize
\begin{tabularx}{\columnwidth}{@{}L{4.2cm}X@{}}
\toprule
\textbf{Attempt} & \textbf{Result} \\
\midrule
Disable the vendor extension & no change \\
Disable subgroup-partitioned ops & no change \\
Force single-threaded pipeline compilation & no change \\
Disable the pipeline library & no change \\
Disable pipeline caching entirely & no change \\
Force synchronous shader compilation & no change \\
Disable descriptor-buffer usage & no change \\
Disable ray tracing & no change \\
Disable mesh shaders & no change \\
Force a conservative memory model & no change \\
Reduce shader-compiler thread count & no change \\
Substitute the other vendor's modules & \textbf{9 $\to$ 4 dialogs, GPU to
                                        100\%} \\
\bottomrule
\end{tabularx}
\end{table}

Eleven configuration changes with no effect and one with a large effect is itself
a signal. When a component's entire exposed configuration surface fails to move
the failure, the component is probably not the one to be configuring.

\subsection{A verification script that was itself wrong}

The override mechanism prints a line for each substituted module, so the harness
checked for those lines to confirm the override had taken effect. It reported
success on runs where nothing had been substituted, because of the \verb|grep -c|
defect described in Section~\ref{sec:deadlock}. Two sweep results had to be
discarded and re-run.

\begin{quote}
\textbf{Every instrument in this project that was not deliberately tested against
a known-bad input was, at some point, wrong.} The harness, the override
verifier, the shader-count comparison, and the process matcher each produced a
confident false reading before being corrected.
\end{quote}

At this point the shader hypothesis was exhausted: a real and reproducible
improvement, no clean explanation, and no path to a shippable configuration. The
next section changes a variable that had been held constant since the beginning
of the project, and the problem disappears entirely.
